\cotefile{84}
\scanpage{51}
% manuscript line(s) 1
\[
N=(T^2\deleted{\unsure{+}}\unsure{-}\delta)/4,
\]
% manuscript line(s) 2–3
c'est une section de $\Gamma\,\mathrm{Sym}^2(E)$, grâce à la condition $T_0^2\equiv\delta\ (4)$. Vérifions que \unsure{le|la} \deleted{fonction} pol.
% manuscript line(s) 4
\[
Q(x)=x^2-Tx+N
\]
% manuscript line(s) 5–6
sur $P_1$ est à valeurs \deleted{sur} $L^{\otimes2}\subset\mathrm{Sym}^2(E)$. Il suffit de prouver qu'il en est ainsi de
% manuscript line(s) 7
\[
4Q(x)=Q(2x)
\]
% manuscript line(s) 8
et on peut poser
% manuscript line(s) 9
\[
2x=T+z
\]
% manuscript line(s) 10
on trouve
% manuscript line(s) 11
\[
Q(T+z)=T^2+2zT+z^2-\underset{\substack{\shortparallel\\\psi(T+z)}}{2}T(T+z)+\underset{\substack{\shortparallel\\\psi(T+z)^2}}{4}N\qquad\text{i.e.}
\]
% manuscript line(s) 12
\[
\boxed{Q(T+z)=\deleted{\unsure{4T^2+}}\,z^2\deleted{\unsure{+}}\unsure{-}\delta}
\]
% manuscript line(s) 13
et on a bien $z^2\deleted{\unsure{+}}\unsure{-}\delta\in\Gamma L^{\otimes2}$ \quad OK.

% manuscript line(s) 14
Description de $\mathcal A$. \unsure{On|Or|Car|On a},
% manuscript line(s) 15
\[
\mathcal A\longrightarrow(\underline O\times L^\vee)\qquad x\longrightarrow(\operatorname{Tr}(x),\varphi(x))
\]
% manuscript line(s) 16–17
est injectif, car le noyau est formé des $x\in\Gamma(\underset{\substack{\cap\\\mathcal A}}{\underline O})$ tels que $\operatorname{Tr}x=0$ i.e. $2x=0$ i.e. $x=0$.
% manuscript line(s) 18
Donc
% manuscript line(s) 18
\[
\mathcal A\subset\underline O\times L^\vee
\]
% manuscript line(s) 19
On trouve
% manuscript line(s) 19
\[
\Gamma\mathcal A\simeq\left\{(b,\unsure{r|\gamma})\in\Gamma\underline O\times\Gamma L^\vee\ \left|\ \boxed{\deleted{\unsure{r_0}}b_0=\deleted{\unsure{b_0}}r_0T_0}\right.\right\}
\]
% manuscript line(s) 20
i.e. $b\in\Gamma\underline O_X$ relève $\underset{\substack{\in\\\Gamma L_0^\vee}}{\deleted{\unsure{b_0}}r_0}\cdot\underset{\substack{\in\\\Gamma L_0}}{T_0}\in\Gamma\underline O_{X_0}$.
% manuscript line(s) 21–22
\diagram{\begin{tikzcd}[ampersand replacement=\&,column sep=large] \unsure{0}\arrow[r]\&{\underline O_X}^{\unsure{0}}\arrow[r,"\xi\mapsto2\xi"]\&\mathcal A\arrow[r,"{(b,r)\mapsto r}","\mathrm{surj}"']\&L^\vee\arrow[r,dash]\&0\end{tikzcd}}
% manuscript line(s) 23
Ainsi l'élém\textsuperscript{t} $\uline{1}$ de $\mathcal A$ est décrit par $(2,0)$
% archive sheet number: 50
